\section{Dynamic Organizational Choice, Bellman System, and Value-Function Existence}

This phase studies the fixed-price dynamic block of the benchmark model. Throughout this section, $(w,p_m,M)$ are taken as given. Therefore the results below establish existence, uniqueness, policy existence, and numerical solvability of the Bellman system conditional on fixed prices. They are not yet a proof of full stationary competitive equilibrium existence.

\subsection{Benchmark Fidelity Check}

\paragraph{Step 1: exact benchmark objects relevant for the present task.}
The authoritative sources for this phase are the benchmark manuscript, the benchmark PDF, and the project dictionary. Cross-checking those sources, the exact fixed-price objects relevant for the present dynamic block are:
\begin{itemize}
    \item $V_A(z;M)$, $V_B(z;M)$, and $V_C(z;M)$: endogenous dynamic value functions for incumbent firms. These are dynamic, endogenous, and structural-only objects.
    \item $\Pi_A(z;w,M)$, $\Pi_B(z;p_m,M)$, and $\Pi_C(z;p_m,w,M)$: optimized static operating payoffs imported from Phase 2. These are derived static objects. They are not new primitives introduced in Phase 3.
    \item $Q_A(\cdot\mid z,M)$, $Q_B(\cdot\mid z,M)$, and $Q_C(\cdot\mid z,M)$: continuation kernels for next period's capability state conditional on the organizational form chosen this period. These are primitive transition objects for the dynamic block.
    \item $\kappa_{sj}(M)$: switching cost from current organizational type $s$ to current-period organizational type $j$. This is a primitive dynamic cost object.
    \item $W_{sj}(z;M)$: derived choice-specific value for a current type-$s$ incumbent that chooses current organization $j$.
    \item $\mathcal Z_{sj}(M)$: derived benchmark choice region, defined as the set of states in which action $j$ solves the Bellman problem for current type $s$.
\end{itemize}

\paragraph{Step 2: exact formulas, argument lists, and interpretation.}
The exact benchmark static payoff objects that enter the Bellman system are
\[
\Pi_A(z;w,M),\quad \Pi_B(z;p_m,M),\quad \Pi_C(z;p_m,w,M).
\]
They must keep these exact argument lists. We do not rewrite them as $\Pi_j(z;w,p_m,M)$. The dynamic endogenous objects to be solved in this phase are
\[
V_A(z;M),\quad V_B(z;M),\quad V_C(z;M).
\]
Economically, the $\Pi_j$ objects summarize the best within-period operating payoff available to organizational form $j$ at state $z$ under fixed prices and policy regime $M$, while the $V_j$ objects summarize the total present discounted value of beginning the period as an incumbent of type $j$ before making the current-period organizational decision.

\paragraph{Step 3: whether the present task can be completed faithfully from the benchmark as written.}
Yes. Holding $(w,p_m,M)$ fixed, the benchmark contains enough structure to:
\begin{enumerate}
    \item construct the incumbent dynamic problem directly from the benchmark timing;
    \item define the Bellman equations for $V_A$, $V_B$, and $V_C$ without changing any benchmark object;
    \item place the Bellman system on an explicit function space;
    \item define a Bellman operator and verify the conditions of Banach's Fixed Point Theorem; and
    \item describe the benchmark argmax correspondence, together with the extra convention required if one wants a single-valued measurable selector.
\end{enumerate}
This means the fixed-price Bellman block can be derived faithfully. What is \emph{not} shown in this phase is full stationary equilibrium existence, because prices are taken as given throughout the present section.

\paragraph{Step 4: inconsistency or ambiguity that must be flagged before proceeding.}
Two benchmark-fidelity issues must be stated explicitly before any derivation begins.
First, there is an argument-list tension across benchmark sources. The dynamic shorthand in the manuscript sometimes writes a generic object of the form $\Pi_j(z;w,p_m,M)$, while the canonical static objects remain
\[
\Pi_A(z;w,M),\qquad \Pi_B(z;p_m,M),\qquad \Pi_C(z;p_m,w,M).
\]
These are not literally identical notational conventions. In this phase we preserve benchmark fidelity by keeping the exact object-specific payoff definitions and treating the generic notation only as shorthand for the appropriate type-$j$ optimized payoff.

Second, there is a tie-resolution tension. The benchmark describes organizational choice through the states in which each action solves the Bellman problem. If ties occur, then the benchmark solving regions can overlap. Therefore a single-valued policy function or a row-stochastic transition matrix requires either an additional tie-breaking convention or an additional uniqueness condition. We do not silently impose either one inside the Bellman existence proof.

\subsection{Part 1: Dynamic Problem, Timing, and Choice-Specific Values}

\paragraph{Benchmark object restatement.}
Fix an incumbent firm with inherited organizational type $s\in\{A,B,C\}$ and inherited capability state $z\in Z$. The firm's current-period feasible action set is
\[
\mathcal J(s)=\{A,B,C,\mathrm{exit}\}.
\]
For later shorthand, define
\[
\mathcal A:=\{A,B,C,\mathrm{exit}\}.
\]
When we write $V_s(z;M)$ with $s\in\{A,B,C\}$, we mean the corresponding coordinate of the value-function triple:
\[
V_s(z;M)=
\begin{cases}
V_A(z;M), & s=A,\\
V_B(z;M), & s=B,\\
V_C(z;M), & s=C.
\end{cases}
\]
This is only shorthand. The benchmark endogenous objects remain the three named functions $V_A$, $V_B$, and $V_C$.

\paragraph{Economic role.}
This block describes how an incumbent converts its current organizational identity and current capability into a current organizational decision. The decision is not merely about how much to produce or how much to research. It is about which organizational form the firm will operate under during the current period. That choice matters immediately because it determines which static operating problem from Phase 2 is available in the same period. It also matters dynamically because the chosen organization determines the continuation kernel for next period's capability state and therefore changes the continuation value.

\paragraph{Mathematical problem.}
Fix a current type $s\in\{A,B,C\}$, a state $z\in Z$, fixed prices $(w,p_m)$, and a fixed regime $M$. The incumbent chooses a current-period organization $j\in\mathcal J(s)$ to solve
\[
V_s(z;M)=\max_{j\in\mathcal J(s)}W_{sj}(z;M).
\]
The left-hand side is the maximal attainable total value from starting the period as current type $s$ at state $z$. The objects inside the maximum, $W_{sj}(z;M)$, are the values attached to each feasible current-period organizational action.

\paragraph{Admissible set.}
For every current type $s$, the feasible set $\mathcal J(s)$ has exactly four elements:
\[
\mathcal J(s)=\{A,B,C,\mathrm{exit}\}.
\]
Therefore the feasible set is nonempty, finite, and independent of $z$ in the benchmark. This point matters mathematically. Because the set of feasible current-period organizational actions is finite, we do not need compact-choice-set machinery to guarantee that the maximum is attained. Once each choice-specific value is a finite real number, the maximum over the four feasible actions exists automatically.

\paragraph{Timing written slowly.}
The benchmark timing can be written as the following sequence.
\begin{enumerate}
    \item The firm enters the period with inherited organization $s$ and inherited capability state $z$.
    \item The firm chooses its \emph{current-period} organizational action $j\in\mathcal J(s)$.
    \item If $j\neq s$, the firm immediately pays the current-period switching cost $\kappa_{sj}(M)$. If $j=s$, there is no reorganization, so we interpret $\kappa_{ss}(M)=0$.
    \item Conditional on choosing an active type $j\in\{A,B,C\}$, the firm solves the within-period operating problem associated with that chosen organization. Phase 2 already solved those within-period problems and summarized them by the optimized current payoff objects $\Pi_A$, $\Pi_B$, and $\Pi_C$.
    \item If the firm chooses exit, it receives the benchmark exit payoff normalized to zero and obtains no continuation value inside this incumbent problem.
    \item If the firm chooses an active type $j\in\{A,B,C\}$, the next-period capability state $z'$ is drawn from the type-specific continuation kernel $Q_j(\cdot\mid z,M)$, and next period begins with inherited organization $j$.
\end{enumerate}

\paragraph{Why the current-period organizational choice must come before the current-period operating controls.}
This point deserves a deliberately slow explanation. In the benchmark, the current operating payoff is not a generic object independent of organizational form. Instead, the current operating block depends on whether the firm operates as type $A$, type $B$, or type $C$ during the current period. A type-$A$ firm solves the integrated-firm problem from Phase 2 and therefore earns $\Pi_A(z;w,M)$. A type-$B$ firm solves the research-specialized problem from Phase 2 and therefore earns $\Pi_B(z;p_m,M)$. A type-$C$ firm solves the manufacturing-specialized problem from Phase 2 and therefore earns $\Pi_C(z;p_m,w,M)$. These are different optimized same-period payoff objects because the benchmark assigns different same-period feasibility sets and different same-period commercialization possibilities to the three organizations. Therefore the model must know the current-period organization before it can even identify which within-period operating problem is relevant. Organization is the outer dynamic choice. Current operating controls are inner static choices already summarized by $\Pi_A$, $\Pi_B$, and $\Pi_C$.

\paragraph{Choice-specific values defined one by one.}
For inherited type $s\in\{A,B,C\}$, define the benchmark choice-specific values
\begin{align*}
W_{sA}(z;M)
&=
\Pi_A(z;w,M)-\kappa_{sA}(M)
+\beta\int_Z V_A(z';M)\,Q_A(dz'\mid z,M),\\
W_{sB}(z;M)
&=
\Pi_B(z;p_m,M)-\kappa_{sB}(M)
+\beta\int_Z V_B(z';M)\,Q_B(dz'\mid z,M),\\
W_{sC}(z;M)
&=
\Pi_C(z;p_m,w,M)-\kappa_{sC}(M)
+\beta\int_Z V_C(z';M)\,Q_C(dz'\mid z,M),\\
W_{s,\mathrm{exit}}(z;M)
&=0.
\end{align*}

\paragraph{Economic meaning of each term in $W_{sj}(z;M)$.}
Each active choice-specific value $W_{sj}(z;M)$ has four ingredients.
\begin{enumerate}
    \item The first term, $\Pi_j$, is the optimized current-period operating payoff obtained by functioning as organization $j$ during the current period.
    \item The second term, $-\kappa_{sj}(M)$, is the one-time current-period cost of reorganizing from current type $s$ into current type $j$.
    \item The third ingredient is the discount factor $\beta$, which converts next period's value into current units.
    \item The fourth ingredient is the expected continuation value, computed using the transition kernel associated with the organization chosen this period.
\end{enumerate}
The exit value is written separately as $W_{s,\mathrm{exit}}(z;M)=0$. The normalization means that once the firm exits this industry problem, it receives neither a positive continuation value nor a negative continuation liability inside the incumbent Bellman block. The important point is that the current organizational choice compares four total values: operate this period as $A$, operate this period as $B$, operate this period as $C$, or leave the incumbent problem entirely.

\subsection{Part 2: Bellman Equations}

For each inherited type $s\in\{A,B,C\}$, the Bellman problem is
\[
V_s(z;M)=\max_{j\in\mathcal J(s)}W_{sj}(z;M).
\]
Since $\mathcal J(s)=\{A,B,C,\mathrm{exit}\}$ for every current type, the benchmark Bellman system is
\begin{align*}
V_A(z;M)&=\max\{W_{AA}(z;M),W_{AB}(z;M),W_{AC}(z;M),0\},\\
V_B(z;M)&=\max\{W_{BA}(z;M),W_{BB}(z;M),W_{BC}(z;M),0\},\\
V_C(z;M)&=\max\{W_{CA}(z;M),W_{CB}(z;M),W_{CC}(z;M),0\}.
\end{align*}
Substituting the definitions of the $W_{sj}$ terms yields the fully expanded system:
\begin{align*}
V_A(z;M)
&=
\max\Bigl\{
\Pi_A(z;w,M)-\kappa_{AA}(M)+\beta\!\int_Z\!V_A(z';M)Q_A(dz'\mid z,M),\\
&\hspace{1.75cm}
\Pi_B(z;p_m,M)-\kappa_{AB}(M)+\beta\!\int_Z\!V_B(z';M)Q_B(dz'\mid z,M),\\
&\hspace{1.75cm}
\Pi_C(z;p_m,w,M)-\kappa_{AC}(M)+\beta\!\int_Z\!V_C(z';M)Q_C(dz'\mid z,M),\,0
\Bigr\},\\
V_B(z;M)
&=
\max\Bigl\{
\Pi_A(z;w,M)-\kappa_{BA}(M)+\beta\!\int_Z\!V_A(z';M)Q_A(dz'\mid z,M),\\
&\hspace{1.75cm}
\Pi_B(z;p_m,M)-\kappa_{BB}(M)+\beta\!\int_Z\!V_B(z';M)Q_B(dz'\mid z,M),\\
&\hspace{1.75cm}
\Pi_C(z;p_m,w,M)-\kappa_{BC}(M)+\beta\!\int_Z\!V_C(z';M)Q_C(dz'\mid z,M),\,0
\Bigr\},\\
V_C(z;M)
&=
\max\Bigl\{
\Pi_A(z;w,M)-\kappa_{CA}(M)+\beta\!\int_Z\!V_A(z';M)Q_A(dz'\mid z,M),\\
&\hspace{1.75cm}
\Pi_B(z;p_m,M)-\kappa_{CB}(M)+\beta\!\int_Z\!V_B(z';M)Q_B(dz'\mid z,M),\\
&\hspace{1.75cm}
\Pi_C(z;p_m,w,M)-\kappa_{CC}(M)+\beta\!\int_Z\!V_C(z';M)Q_C(dz'\mid z,M),\,0
\Bigr\}.
\end{align*}
The problem is recursive because the unknown object we seek today, namely the value function, reappears inside tomorrow's continuation term.

\subsection{Part 3: Function Space, Norm, and Bellman Operator}

From Phase 1, let $Z$ be a nonempty compact metric space and let $\mathcal B(Z)$ denote its Borel sigma algebra. Define
\[
\mathcal C_b(Z)=\{f:Z\to\mathbb R\mid f\text{ is bounded and continuous}\},
\qquad
\|f\|_\infty=\sup_{z\in Z}|f(z)|.
\]
Because there are three active organizational types, the candidate value-function space is
\[
\mathcal V:=\mathcal C_b(Z)^3.
\]
For a generic triple $U=(U_A,U_B,U_C)\in\mathcal V$, use the product sup norm
\[
\|U\|:=\max\{\|U_A\|_\infty,\|U_B\|_\infty,\|U_C\|_\infty\}.
\]

To define one Bellman operator compactly, introduce auxiliary proof notation. For $j\in\{A,B,C\}$, define
\[
(\mathcal Q_jU)(z):=\int_Z U_j(z')Q_j(dz'\mid z,M),
\qquad
(\mathcal Q_{\mathrm{exit}}U)(z):=0.
\]
For inherited type $s\in\{A,B,C\}$ and action $a\in\{A,B,C,\mathrm{exit}\}$, define
\[
\mathcal W_{sa}(z;M;U):=
\begin{cases}
\Pi_A(z;w,M)-\kappa_{sA}(M)+\beta(\mathcal Q_AU)(z), & a=A,\\
\Pi_B(z;p_m,M)-\kappa_{sB}(M)+\beta(\mathcal Q_BU)(z), & a=B,\\
\Pi_C(z;p_m,w,M)-\kappa_{sC}(M)+\beta(\mathcal Q_CU)(z), & a=C,\\
0, & a=\mathrm{exit}.
\end{cases}
\]
At the fixed point $V^*$, the benchmark choice-specific values are recovered by
\[
W_{sa}(z;M)=\mathcal W_{sa}(z;M;V^*).
\]
The Bellman operator is
\[
(\mathcal T_MU)_s(z):=\max_{a\in\{A,B,C,\mathrm{exit}\}}\mathcal W_{sa}(z;M;U),
\qquad s\in\{A,B,C\}.
\]
The fixed-price Bellman problem is therefore equivalent to the fixed-point equation $V=\mathcal T_MV$.

\subsection{Assumptions}

\paragraph{Benchmark ingredients carried into the proof.}
The Bellman proof begins from benchmark objects that already exist before Phase 3 starts. From Phase 1 we inherit the state space $Z$ as a nonempty compact metric space equipped with its Borel sigma algebra. From Phase 2 we inherit the three optimized static payoff objects $\Pi_A(z;w,M)$, $\Pi_B(z;p_m,M)$, and $\Pi_C(z;p_m,w,M)$. From the benchmark dynamic block we inherit the continuation kernels $Q_A(\cdot\mid z,M)$, $Q_B(\cdot\mid z,M)$, and $Q_C(\cdot\mid z,M)$, the switching-cost object $\kappa_{sj}(M)$, and the discount factor $\beta$. The assumptions listed below therefore do not redesign the model. Their role is to make explicit which regularity properties of those already-defined benchmark objects are needed to prove Bellman existence, uniqueness, and policy measurability.

\paragraph{How to read R1--R4.}
R1, R3, and R4 are direct regularity statements about benchmark objects already used in the model. R2 deserves separate emphasis. The benchmark defines the continuation kernels $Q_j$, but the fixed-point proof additionally needs the Feller property so that continuation integration preserves continuity of candidate value functions. We therefore state R2 openly as a supplemental regularity assumption used by the proof. This is exactly the kind of assumption that should be explicit in a technical appendix rather than left implicit inside the phrase ``standard dynamic programming arguments.''

\begin{assumption}[R1: bounded-continuous flow objects]
The optimized payoffs satisfy $\Pi_A(\cdot;w,M)\in\mathcal C_b(Z)$, $\Pi_B(\cdot;p_m,M)\in\mathcal C_b(Z)$, and $\Pi_C(\cdot;p_m,w,M)\in\mathcal C_b(Z)$.
\end{assumption}

\begin{assumption}[R2: Feller probability kernels]
For each $j\in\{A,B,C\}$ and each $(z,M)$, $Q_j(\cdot\mid z,M)$ is a Borel probability measure on $Z$. Moreover, for each $f\in\mathcal C_b(Z)$,
\[
z\mapsto\int_Z f(z')Q_j(dz'\mid z,M)
\]
is continuous on $Z$.
\end{assumption}

\begin{assumption}[R3: discounting]
$\beta\in(0,1)$.
\end{assumption}

\begin{assumption}[R4: finite switching costs]
$|\kappa_{sj}(M)|<\infty$ for all $s,j\in\{A,B,C\}$.
\end{assumption}

\paragraph{A theorem used inside the completeness proof.}
One theorem will be used explicitly below: the Uniform Limit Theorem. It states that if $X$ is a metric space, each $f^n:X\to\mathbb R$ is continuous, and $f^n\to f$ uniformly on $X$, then $f$ is continuous on $X$. We record it here so that the proof can cite an exact theorem by name rather than appealing to a vague continuity argument. In our application, the domain is the compact metric space $Z$, the approximating functions belong to $\mathcal C_b(Z)$, and uniform convergence will be obtained from the Cauchy property under the sup norm.

\subsection{Part 4: Well-Posedness of the Bellman System}

We now prove that the fixed-price Bellman block is mathematically well posed. Because the Bellman equations are recursive equations in unknown functions, well-posedness here means more than simply writing those equations down. We must show that there is a function space in which the Bellman operator is well defined, that the operator maps that space into itself, that it shrinks distances under an explicit norm, and therefore that the recursive system has a unique solution. Throughout this subsection, prices $(w,p_m)$ and policy regime $M$ remain fixed. This is still a fixed-price dynamic block, not yet a proof of stationary equilibrium existence.

\paragraph{What assumptions are being used here, and why?}
Assumptions R1--R4 combine benchmark objects with the extra regularity conditions needed for the proof. R1 records that the optimized Phase 2 payoff objects are bounded and continuous. R2 records that the continuation kernels are Borel probability measures and satisfy the Feller continuity property needed to preserve continuity under integration. R3 is the discounting condition $\beta\in(0,1)$ that later gives the contraction modulus. R4 records finiteness of switching costs so that current-period reorganization does not generate undefined values. None of these assumptions changes the benchmark economic problem; they only make explicit the regularity conditions needed for a full fixed-point proof.

\paragraph{A theorem used inside the completeness proof.}
We will use the Uniform Limit Theorem. It states that if $X$ is a metric space, each function $f^n:X\to\mathbb R$ is continuous, and $f^n\to f$ uniformly on $X$, then the limit function $f$ is continuous on $X$. The assumptions therefore have to be checked in exactly that order: metric-space domain, continuity of every approximating function, and uniform rather than merely pointwise convergence.

\begin{lemma}
Under R1--R4, $(\mathcal C_b(Z),\|\cdot\|_\infty)$ is complete. Consequently, $(\mathcal V,\|\cdot\|)$ is complete.
\end{lemma}

\begin{proof}
Take a Cauchy sequence $\{f^n\}$ in $(\mathcal C_b(Z),\|\cdot\|_\infty)$. By definition of the sup norm, this means that for every $\varepsilon>0$ there exists $N(\varepsilon)$ such that
\[
\|f^n-f^m\|_\infty<\varepsilon
\qquad\text{for all }n,m\ge N(\varepsilon).
\]

\medskip\noindent\textbf{Step 1: construct a pointwise candidate limit.}
Fix any state $z\in Z$. Then
\[
|f^n(z)-f^m(z)|\le \|f^n-f^m\|_\infty<\varepsilon
\qquad\text{for all }n,m\ge N(\varepsilon).
\]
So $\{f^n(z)\}$ is Cauchy in $\mathbb R$, hence converges to some limit $f(z)\in\mathbb R$.

\medskip\noindent\textbf{Step 2: upgrade pointwise convergence to uniform convergence.}
We now prove that the convergence is uniform, because continuity of the limit will later rely on that stronger mode of convergence. Fix $n\ge N(\varepsilon)$. For all $m\ge N(\varepsilon)$,
\[
\sup_{z\in Z}|f^n(z)-f^m(z)|<\varepsilon.
\]
Let $m\to\infty$. Since $f^m(z)\to f(z)$ pointwise,
\[
|f^n(z)-f(z)|\le\varepsilon\qquad\text{for every }z\in Z.
\]
Taking the supremum over $z$ gives
\[
\|f^n-f\|_\infty\le \varepsilon.
\]
So $f^n\to f$ uniformly.

\medskip\noindent\textbf{Step 3: prove that the limit function belongs to $\mathcal C_b(Z)$.}
Each $f^n$ belongs to $\mathcal C_b(Z)$, so each $f^n$ is continuous. The domain $Z$ is a metric space, and we have just shown that $f^n\to f$ uniformly. Therefore all assumptions of the Uniform Limit Theorem are satisfied, so $f$ is continuous on $Z$.

It remains to show boundedness. Because the sequence is Cauchy in sup norm, choose $N$ such that $\|f^n-f^N\|_\infty\le 1$ for all $n\ge N$. Passing to the limit gives $\|f-f^N\|_\infty\le 1$, hence
\[
\|f\|_\infty\le \|f-f^N\|_\infty+\|f^N\|_\infty\le 1+\|f^N\|_\infty<\infty.
\]
Therefore $f\in\mathcal C_b(Z)$, so $\mathcal C_b(Z)$ is complete.

\medskip\noindent\textbf{Step 4: pass from one coordinate to the product space.}
Now let $U^n=(U_A^n,U_B^n,U_C^n)$ be a Cauchy sequence in $\mathcal V$. Because the product norm is the maximum of the three coordinate sup norms, each coordinate sequence is Cauchy in $\mathcal C_b(Z)$:
\[
\|U_A^n-U_A^m\|_\infty\le \|U^n-U^m\|,
\quad
\|U_B^n-U_B^m\|_\infty\le \|U^n-U^m\|,
\quad
\|U_C^n-U_C^m\|_\infty\le \|U^n-U^m\|.
\]
Each coordinate therefore converges in $\mathcal C_b(Z)$, and the three coordinate limits define a limit in $\mathcal V$. Hence $(\mathcal V,\|\cdot\|)$ is complete.
\end{proof}

\begin{lemma}
Under R1--R4, for every $U\in\mathcal V$, $\mathcal T_MU\in\mathcal V$.
\end{lemma}

\begin{proof}
Fix $U\in\mathcal V$ and inherited type $s\in\{A,B,C\}$. We must show that $(\mathcal T_MU)_s$ is both bounded and continuous on $Z$.

\medskip\noindent\textbf{Step 1: each continuation term is well defined and bounded.}
For each active action $j\in\{A,B,C\}$,
\begin{align*}
|(\mathcal Q_jU)(z)|
&=
\left|\int_Z U_j(z')Q_j(dz'\mid z,M)\right|\\
&\le
\int_Z |U_j(z')|Q_j(dz'\mid z,M)\\
&\le
\int_Z \|U_j\|_\infty Q_j(dz'\mid z,M)\\
&= 
\|U_j\|_\infty.
\end{align*}
So each continuation term is a finite real number for every $z\in Z$. The finiteness is not automatic unless we use both inputs: $U_j$ is bounded because $U\in\mathcal V$, and $Q_j(\cdot\mid z,M)$ has total mass one because it is a probability measure.

\medskip\noindent\textbf{Step 2: each active choice-specific index is bounded.}
Define the finite constants
\[
\overline \Pi
:=
\max\left\{
\sup_{z\in Z}|\Pi_A(z;w,M)|,\,
\sup_{z\in Z}|\Pi_B(z;p_m,M)|,\,
\sup_{z\in Z}|\Pi_C(z;p_m,w,M)|
\right\},
\]
and
\[
\overline \kappa
:=
\max_{s,j\in\{A,B,C\}}|\kappa_{sj}(M)|.
\]
These constants are finite by R1 and R4. Consider, for example, the action-$A$ index:
\begin{align*}
|\mathcal W_{sA}(z;M;U)|
&=
\left|\Pi_A(z;w,M)-\kappa_{sA}(M)+\beta(\mathcal Q_AU)(z)\right|\\
&\le
|\Pi_A(z;w,M)|+|\kappa_{sA}(M)|+\beta|(\mathcal Q_AU)(z)|\\
&\le
\overline \Pi+\overline \kappa+\beta\|U_A\|_\infty\\
&\le
\overline \Pi+\overline \kappa+\beta\|U\|.
\end{align*}
The same calculation applies to the action-$B$ and action-$C$ indices. For the exit action, $\mathcal W_{s,\mathrm{exit}}(\cdot;M;U)\equiv 0$, which is trivially bounded.

\medskip\noindent\textbf{Step 3: each active choice-specific index is continuous.}
By R2, if $f\in\mathcal C_b(Z)$, then $z\mapsto \int_Z f(z')Q_j(dz'\mid z,M)$ is continuous. Because $U_j\in\mathcal C_b(Z)$ for each $j$, it follows that each continuation term $(\mathcal Q_jU)(\cdot)$ is continuous. The payoff term $\Pi_j(\cdot)$ is continuous by R1, the switching cost $\kappa_{sj}(M)$ is a constant with respect to $z$, and sums of continuous functions remain continuous. Therefore, for active actions,
\[
\mathcal W_{sj}(\cdot;M;U)\in\mathcal C_b(Z).
\]

\medskip\noindent\textbf{Step 4: take the maximum across actions.}
The function $(\mathcal T_MU)_s$ is the maximum of four bounded continuous functions: the three active indices and the exit index. It is useful to say one extra sentence here instead of burying the argument inside the word ``maximum.'' For any two real numbers $x$ and $y$,
\[
\max\{x,y\}=\frac{x+y+|x-y|}{2}.
\]
Therefore, if $f$ and $g$ are continuous functions, then $\max\{f,g\}$ is continuous because addition, subtraction, scalar multiplication, and absolute value preserve continuity. Repeating that argument finitely many times shows that the maximum of finitely many continuous functions is continuous. Boundedness is immediate because the maximum of finitely many bounded functions is bounded. Hence $(\mathcal T_MU)_s\in\mathcal C_b(Z)$. Because this argument applies to $s=A$, $s=B$, and $s=C$, we conclude that $\mathcal T_MU\in\mathcal V$.
\end{proof}

\begin{lemma}
Under R1--R4, for all $U,\widehat U\in\mathcal V$,
\[
\|\mathcal T_MU-\mathcal T_M\widehat U\|\le\beta\|U-\widehat U\|.
\]
\end{lemma}

\begin{proof}
Fix inherited type $s$ and state $z$.

\medskip\noindent\textbf{Step 1: control the difference of two finite maxima.}
We first prove the elementary inequality
\[
|\max_i x_i-\max_i y_i|\le\max_i|x_i-y_i|.
\]
Let $i_x$ attain $\max_i x_i$. Then
\[
\max_i x_i-\max_i y_i=x_{i_x}-\max_i y_i\le x_{i_x}-y_{i_x}\le \max_i|x_i-y_i|.
\]
Switching the roles of $x$ and $y$ gives
\[
\max_i y_i-\max_i x_i\le \max_i|x_i-y_i|.
\]
Combining the two inequalities yields the absolute-value result.

\medskip\noindent\textbf{Step 2: compare the Bellman indices action by action.}
Therefore
\[
|(\mathcal T_MU)_s(z)-(\mathcal T_M\widehat U)_s(z)|
\le\max_{a\in\{A,B,C,\mathrm{exit}\}}|\mathcal W_{sa}(z;M;U)-\mathcal W_{sa}(z;M;\widehat U)|.
\]
For the exit action, the difference is zero because the exit value is always normalized to zero, independently of $U$.

For an active action $a=j\in\{A,B,C\}$, note first that the flow payoff $\Pi_j$ and the switching-cost term $\kappa_{sj}(M)$ appear identically in both $\mathcal W_{sj}(z;M;U)$ and $\mathcal W_{sj}(z;M;\widehat U)$. They therefore cancel when the two indices are subtracted. Only the continuation term depends on the candidate value function. Hence
\begin{align*}
|\mathcal W_{sj}(z;M;U)-\mathcal W_{sj}(z;M;\widehat U)|
&=\beta\left|\int_Z \bigl(U_j(z')-\widehat U_j(z')\bigr)Q_j(dz'\mid z,M)\right|\\
&\le\beta\int_Z |U_j(z')-\widehat U_j(z')|Q_j(dz'\mid z,M)\\
&\le\beta\|U_j-\widehat U_j\|_\infty\\
&\le\beta\|U-\widehat U\|.
\end{align*}

\medskip\noindent\textbf{Step 3: pass from one state and one current type to the product norm.}
The previous bound holds for every active action $j$, every inherited type $s$, and every state $z$. Therefore
\[
|(\mathcal T_MU)_s(z)-(\mathcal T_M\widehat U)_s(z)|
\le
\beta\|U-\widehat U\|
\qquad\text{for all }s,z.
\]
Taking the supremum over $z\in Z$ yields
\[
\|(\mathcal T_MU)_s-(\mathcal T_M\widehat U)_s\|_\infty
\le
\beta\|U-\widehat U\|.
\]
Then taking the maximum over inherited types $s\in\{A,B,C\}$ yields
\[
\|\mathcal T_MU-\mathcal T_M\widehat U\|
\le
\beta\|U-\widehat U\|.
\]
So $\mathcal T_M$ is a contraction on $\mathcal V$ with modulus $\beta$.
\end{proof}

\begin{theorem}
Under R1--R4, there exists a unique $V^*=(V_A^*,V_B^*,V_C^*)\in\mathcal V$ such that
\[
\mathcal T_MV^*=V^*.
\]
Equivalently, the fixed-price Bellman system has a unique solution in $\mathcal V$.
\end{theorem}

\begin{proof}
Before invoking any fixed-point theorem, note the exact equivalence between the Bellman system and the fixed-point equation. If $V$ satisfies $V=\mathcal T_MV$, then by definition each coordinate of $\mathcal T_MV$ equals the corresponding Bellman maximization problem, so $V_A$, $V_B$, and $V_C$ solve the Bellman equations. Conversely, if a triple $V$ solves the three Bellman equations, then each coordinate of $V$ equals the corresponding coordinate of $\mathcal T_MV$, so $V=\mathcal T_MV$. Thus solving the Bellman system is exactly the same task as finding a fixed point of $\mathcal T_M$.

We now invoke Banach's Fixed Point Theorem. The theorem states: if $(X,d)$ is a complete metric space and $T:X\to X$ satisfies
\[
d(Tx,Ty)\le q\,d(x,y)\qquad\text{for all }x,y\in X
\]
for some $q\in[0,1)$, then $T$ has a unique fixed point in $X$, and iterates from any starting point converge to that fixed point.

We verify its assumptions one by one.

First, the completeness lemma showed that $(\mathcal V,\|\cdot\|)$ is a complete metric space.

Second, the self-map lemma showed that $\mathcal T_M:\mathcal V\to\mathcal V$.

Third, the contraction lemma showed that
\[
\|\mathcal T_MU-\mathcal T_M\widehat U\|\le \beta\|U-\widehat U\|
\]
for all $U,\widehat U\in\mathcal V$, and R3 gives $\beta\in(0,1)$, hence $\beta<1$.

All assumptions of Banach's theorem are therefore satisfied. The theorem yields a unique fixed point $V^*\in\mathcal V$ with $\mathcal T_MV^*=V^*$. Because the fixed-point equation is exactly the Bellman system, the value-function triple is uniquely determined. The theorem also tells us something constructive, not merely existential: starting from any initial guess in $\mathcal V$, repeated Bellman iteration converges to this same unique fixed point.
\end{proof}

\paragraph{What this theorem has established, and what it has not.}
Part 4 establishes a fixed-price dynamic result. It says that once $(w,p_m,M)$ are held fixed and the regularity assumptions R1--R4 are imposed, the incumbent Bellman system has one and only one value-function triple in $\mathcal V$. That conclusion is already strong enough to support the later policy analysis, because it implies that the continuation values used by incumbents are not ambiguous. At the same time, the theorem has not yet proved that equilibrium prices exist, that entry is well defined, or that stationary distributions exist. Those are later equilibrium questions. The present theorem only solves the incumbent dynamic block conditional on fixed prices and the fixed policy regime.

\subsection{Part 5: Policy Existence, Set-Valued Policies, and Measurable Selectors}

The fixed-point theorem delivers a unique value-function triple, but it does not by itself tell us whether the induced organizational policy is single-valued at every state. That issue has to be handled separately. The benchmark object is the set of actions that solve the Bellman problem, not an already tie-broken selector.

\paragraph{Why policy existence is a separate step from value-function existence.}
Part 4 solved the recursive value equations. Part 5 performs a different logical task. It asks what the solved value functions imply for the firm's actual organizational choices at each state. This distinction is important because uniqueness of the value-function triple does not imply uniqueness of the maximizing action state by state. A unique fixed point can still generate ties inside the Bellman maximization. Therefore policy analysis cannot be skipped as if it were already contained in the fixed-point theorem. The policy object has to be constructed separately from the solved values, and the tie issue has to be handled explicitly.

\paragraph{Benchmark policy objects.}
For each inherited type $s\in\{A,B,C\}$, define the benchmark argmax correspondence
\[
\Gamma_s(z;M):=\arg\max_{a\in\mathcal A}\mathcal W_{sa}(z;M;V^*).
\]
This correspondence tells us which current-period organizational actions solve the Bellman problem for a current type-$s$ firm at state $z$. The benchmark choice-region object is then
\[
\mathcal Z_{sj}(M):=\{z\in Z:j\in\Gamma_s(z;M)\}.
\]
This is exactly the benchmark definition. We are not replacing it with a selector-based alternative.

\begin{proposition}
For each inherited type $s\in\{A,B,C\}$:
\begin{enumerate}
    \item $\Gamma_s(z;M)$ is nonempty for every $z\in Z$;
    \item for each action $j\in\{A,B,C,\mathrm{exit}\}$, the benchmark choice region $\mathcal Z_{sj}(M)$ is a Borel subset of $Z$;
    \item if exactly one action attains the maximum at state $z$, then the policy is single-valued at that state;
    \item if more than one action attains the maximum at state $z$, then the benchmark policy object is set-valued at that state;
    \item if one adds an explicit deterministic tie-breaking order, then one can construct a single-valued measurable selector. That selector is an auxiliary implementation device, not the benchmark primitive object.
\end{enumerate}
\end{proposition}

\begin{proof}
\medskip\noindent\textbf{Step 1: argmax existence.}
Fix a current type $s$ and a state $z$. The Bellman problem compares the four real numbers
\[
\mathcal W_{sA}(z;M;V^*),\quad
\mathcal W_{sB}(z;M;V^*),\quad
\mathcal W_{sC}(z;M;V^*),\quad
0.
\]
These are finite real numbers because $V^*\in\mathcal V$, so the choice-specific values are well defined. A nonempty finite set of real numbers always attains its maximum. Hence $\Gamma_s(z;M)$ is nonempty.

\medskip\noindent\textbf{Step 2: an explicit formula for the benchmark choice regions.}
Fix $s$ and $j$. By definition, $z\in\mathcal Z_{sj}(M)$ if and only if action $j$ attains at least as much value as every competitor $\ell\in\mathcal A$. Therefore
\[
\mathcal Z_{sj}(M)
=
\bigcap_{\ell\in\mathcal A}
\{z\in Z:\mathcal W_{sj}(z;M;V^*)\ge \mathcal W_{s\ell}(z;M;V^*)\}.
\]
This formula rewrites the benchmark choice region as a finite intersection of comparison sets.

\medskip\noindent\textbf{Step 3: measurability of the benchmark choice regions.}
Fix a competitor $\ell\in\mathcal A$. Consider the difference
\[
z\mapsto \mathcal W_{sj}(z;M;V^*)-\mathcal W_{s\ell}(z;M;V^*)
\]
The self-map proof already established continuity of each $\mathcal W_{sj}(\cdot;M;V^*)$, so the difference of two such functions is continuous. Therefore the set
\[
\{z\in Z:\mathcal W_{sj}(z;M;V^*)-\mathcal W_{s\ell}(z;M;V^*)\ge 0\}
\]
is the preimage of the closed set $[0,\infty)$ under a continuous map. Hence it is closed. A finite intersection of closed sets is closed. A closed subset of the metric space $Z$ is Borel measurable. Therefore $\mathcal Z_{sj}(M)$ is a Borel subset of $Z$.

\medskip\noindent\textbf{Step 4: single-valued versus set-valued policy.}
If exactly one action $j$ attains the maximum at state $z$, then the argmax set contains exactly one element:
\[
\Gamma_s(z;M)=\{j\}.
\]
In that case the policy is single-valued at $z$. If two or more actions attain the same maximal value, then $\Gamma_s(z;M)$ contains multiple elements, so the benchmark policy object is set-valued at that state.

\medskip\noindent\textbf{Step 5: measurable selector after adding an explicit tie-breaking convention.}
Suppose we impose the deterministic priority order
\[
A\succ B\succ C\succ \mathrm{exit}.
\]
Define
\begin{align*}
D_{sA}(M)&:=\mathcal Z_{sA}(M),\\
D_{sB}(M)&:=\mathcal Z_{sB}(M)\setminus D_{sA}(M),\\
D_{sC}(M)&:=\mathcal Z_{sC}(M)\setminus\bigl(D_{sA}(M)\cup D_{sB}(M)\bigr),\\
D_{s,\mathrm{exit}}(M)&:=\mathcal Z_{s,\mathrm{exit}}(M)\setminus\bigl(D_{sA}(M)\cup D_{sB}(M)\cup D_{sC}(M)\bigr).
\end{align*}
These sets are Borel because they are formed from Borel sets by finite unions, intersections, and complements. Now define
\[
g_s(z;M)=
\begin{cases}
A, & z\in D_{sA}(M),\\
B, & z\in D_{sB}(M),\\
C, & z\in D_{sC}(M),\\
\mathrm{exit}, & z\in D_{s,\mathrm{exit}}(M).
\end{cases}
\]
Because the original argmax set is nonempty at every state, every state belongs to at least one of the original benchmark regions $\mathcal Z_{sj}(M)$. The sequential subtraction assigns the highest-priority maximizing action at each state. Hence $g_s(z;M)\in\Gamma_s(z;M)$ for every $z\in Z$. The function $g_s$ is measurable because it is finite-valued and each inverse image $g_s^{-1}(\{A\})$, $g_s^{-1}(\{B\})$, $g_s^{-1}(\{C\})$, and $g_s^{-1}(\{\mathrm{exit}\})$ is Borel measurable.

\medskip\noindent\textbf{Which theorem did we use for measurable selection?}
None. Because the action set is finite, we do not need Kuratowski--Ryll-Nardzewski or Berge's Maximum Theorem to obtain a measurable selector. Once an explicit tie-breaking rule is imposed, the selector can be constructed directly from measurable benchmark choice regions. This keeps the logical dependence transparent: the benchmark object is the correspondence, while the selector appears only after an extra convention is added.

\medskip\noindent\textbf{Does the benchmark itself assume deterministic policies?}
Yes, but only in a limited sense. The benchmark does not introduce random utility shocks, logit smoothing, or mixed strategies. However, that does not mean the benchmark automatically gives a single-valued policy function before ties are examined. The benchmark-faithful object remains the argmax correspondence $\Gamma_s$ or, equivalently, the family of benchmark choice regions $\mathcal Z_{sj}(M)$. Therefore the right way to describe the model is: it is deterministic once a maximizing action has been selected, but ties must still be treated explicitly whenever later theory or computation requires one action to be chosen at each state.
\end{proof}

\subsection{Choice regions and transition interface}

For the benchmark transition interface, define the choice regions directly from the Bellman problem:
\[
\mathcal Z_{sj}(M)=\{z\in Z:j\in\Gamma_s(z;M)\},
\qquad s\in\{A,B,C\},\ j\in\{A,B,C,\mathrm{exit}\}.
\]
\paragraph{Why these regions are the bridge to Phase 4.}
Phase 3 ends with solved value functions and induced Bellman policy objects. Phase 4, by contrast, needs objects that move mass across organizational forms and state space. The family of measurable sets $\{\mathcal Z_{sj}(M)\}$ is the bridge between those two phases. Economically, each set records the states in which a current type-$s$ firm chooses organizational form $j$ this period. Mathematically, those sets allow the stationary measures of current incumbents to be pushed through the Bellman policy into transition objects that can later be assembled into a transition matrix and stationary-distribution system.

Given stationary measures $\mu_s$, Phase 4 uses
\[
P_{sj}(M)=\frac{\mu_s(\mathcal Z_{sj}(M))}{\mu_s(Z)}.
\]
Because the choice regions are measurable, these set-based transition objects are well defined whenever $\mu_s(Z)>0$. However, benchmark fidelity requires a further clarification. The regions $\mathcal Z_{sj}(M)$ are not automatically disjoint. If ties occur, two or more actions may solve the Bellman problem at the same state, so two or more benchmark choice regions may overlap. Therefore measurability alone is enough to define the set measures above, but it is not by itself enough to guarantee that the resulting matrix is already row-stochastic. A row-stochastic transition matrix requires either a zero-measure-ties condition or an explicit selector.

If later computation requires a single-valued deterministic policy, an auxiliary selector can be introduced only with an explicit tie-breaking convention. For example, under the order $A\succ B\succ C\succ \mathrm{exit}$, define
\begin{align*}
D_{sA}(M)&:=\mathcal Z_{sA}(M),\\
D_{sB}(M)&:=\mathcal Z_{sB}(M)\setminus D_{sA}(M),\\
D_{sC}(M)&:=\mathcal Z_{sC}(M)\setminus\bigl(D_{sA}(M)\cup D_{sB}(M)\bigr),\\
D_{s,\mathrm{exit}}(M)&:=\mathcal Z_{s,\mathrm{exit}}(M)\setminus\bigl(D_{sA}(M)\cup D_{sB}(M)\cup D_{sC}(M)\bigr).
\end{align*}
Then
\[
g_s(z;M)=
\begin{cases}
A, & z\in D_{sA}(M),\\
B, & z\in D_{sB}(M),\\
C, & z\in D_{sC}(M),\\
\mathrm{exit}, & z\in D_{s,\mathrm{exit}}(M)
\end{cases}
\]
is a Borel measurable selector satisfying $g_s(z;M)\in\Gamma_s(z;M)$ for every $z$. This selector is an auxiliary implementation device only. The benchmark object remains the correspondence $\Gamma_s$ or, equivalently, the choice regions $\mathcal Z_{sj}(M)$. The reason to keep this distinction visible is that Phase 4 will use the solved value functions and the induced policy objects to build the transition matrix and the stationary distributions. If the later phase requires disjoint policy regions, the extra selector convention must be stated there explicitly rather than hidden inside Phase 3 notation.

\paragraph{What Phase 4 will inherit from Phase 3.}
The next phase should therefore take three objects as inputs from the present phase. First, it inherits the unique fixed-price value-function triple $V^*$. Second, it inherits the benchmark argmax correspondences $\Gamma_s$. Third, it inherits the measurable benchmark choice regions $\mathcal Z_{sj}(M)$. Those objects are the correct Bellman outputs for constructing endogenous transition matrices, stationary distributions across organizational forms, and the equilibrium blocks that later endogenize prices. Writing this bridge explicitly matters because it keeps the logical order clear: solve the fixed-price dynamic block first, translate solved values into measurable policy objects second, and only then build the stationary transition structure.

\subsection{How the Bellman System and the Value Functions Are Solved}

The preceding subsections have separately defined the Bellman equations, the Bellman operator, the existence-and-uniqueness conditions, and the policy objects. This subsection consolidates those results into one continuous solution sequence. The purpose is to state clearly how the fixed-price dynamic block moves from the optimized static payoffs to the unique value-function triple and, from there, to the policy objects that will be used in Phase 4.

\paragraph{Step 1: construct the choice-specific values.}
For inherited type $s\in\{A,B,C\}$ and state $z\in Z$, the firm faces four actions:
\[
\mathcal A=\{A,B,C,\mathrm{exit}\}.
\]
If the firm chooses organizational form $A$, its current total value is
\[
W_{sA}(z;M)
=
\Pi_A(z;w,M)-\kappa_{sA}(M)+\beta\int_Z V_A(z';M)Q_A(dz'\mid z,M).
\]
If it chooses $B$ or $C$, the definitions are analogous:
\[
W_{sB}(z;M)
=
\Pi_B(z;p_m,M)-\kappa_{sB}(M)+\beta\int_Z V_B(z';M)Q_B(dz'\mid z,M),
\]
\[
W_{sC}(z;M)
=
\Pi_C(z;p_m,w,M)-\kappa_{sC}(M)+\beta\int_Z V_C(z';M)Q_C(dz'\mid z,M).
\]
If it exits, then
\[
W_{s,\mathrm{exit}}(z;M)=0.
\]
Accordingly, the starting point of the dynamic solution is not to guess a closed form for $V_s$. The starting point is to write down the full current value of each currently feasible action. Once that is done, the Bellman problem compares a finite collection of explicitly defined objects.

\paragraph{Step 2: write the value function as the upper envelope of the choice-specific values.}
Once the four choice-specific values are available, the inherited-type value function is simply
\[
V_s(z;M)=\max_{j\in\mathcal A}W_{sj}(z;M).
\]
Equivalently,
\[
V_A(z;M)=\max\{W_{AA}(z;M),W_{AB}(z;M),W_{AC}(z;M),0\},
\]
\[
V_B(z;M)=\max\{W_{BA}(z;M),W_{BB}(z;M),W_{BC}(z;M),0\},
\]
\[
V_C(z;M)=\max\{W_{CA}(z;M),W_{CB}(z;M),W_{CC}(z;M),0\}.
\]
Therefore, in this benchmark, ``solving the Bellman system'' and ``solving for the three value functions $V_A,V_B,V_C$'' are the same task. Once those three value functions are known, the ranking of the choice-specific values, the argmax correspondences, and the choice regions are all pinned down.

\paragraph{Step 3: rewrite the Bellman system as a fixed-point problem.}
The difficulty is that the current unknown value function reappears inside the continuation term. For example,
\[
W_{sA}(z;M)
=
\Pi_A(z;w,M)-\kappa_{sA}(M)+\beta\int_Z V_A(z';M)Q_A(dz'\mid z,M)
\]
still contains the unknown function $V_A$ on the right-hand side. So this is not a finite-dimensional algebraic system. It is a system of functional equations. More precisely, we are looking for a function triple
\[
V=(V_A,V_B,V_C)
\]
such that
\[
V=\mathcal T_MV.
\]
Thus, mathematically, solving the Bellman system is equivalent to finding a fixed point of the Bellman operator. This step turns a recursive system of functional equations into one unified fixed-point problem on a function space.

\paragraph{Step 4: use Banach's theorem to obtain the unique value function.}
To prove that this fixed point is well defined, we place candidate value functions in the function space
\[
\mathcal V=\mathcal C_b(Z)^3
\]
and equip that space with the product sup norm. We then prove three facts:
\begin{enumerate}
    \item $(\mathcal V,\|\cdot\|)$ is complete;
    \item the Bellman operator $\mathcal T_M$ maps $\mathcal V$ into itself;
    \item the Bellman operator satisfies the contraction inequality
    \[
    \|\mathcal T_MU-\mathcal T_M\widehat U\|\le \beta\|U-\widehat U\|.
    \]
\end{enumerate}
Because $\beta\in(0,1)$, $\mathcal T_M$ is a contraction. Banach's Fixed Point Theorem therefore yields a unique fixed point
\[
V^*=(V_A^*,V_B^*,V_C^*)\in\mathcal V
\]
satisfying
\[
\mathcal T_MV^*=V^*.
\]
This is the theoretical solution of the value functions. In other words, the fixed-price Bellman block is not merely well defined; it is uniquely pinned down by the benchmark primitives and the stated regularity conditions.

\paragraph{Step 5: compute that fixed point constructively by value iteration.}
Banach's theorem is useful not only because it gives existence and uniqueness, but also because it gives a constructive algorithm. Starting from any bounded initial guess $V^{(0)}$, repeated iteration
\[
V^{(k+1)}=\mathcal T_MV^{(k)}
\]
converges to the unique fixed point $V^*$. On a discrete grid, this consists of three operations:
\begin{enumerate}
    \item compute continuation values using the transition matrices;
    \item compute the four choice-specific values at each node;
    \item take the maximum across the four actions to obtain the next value-function iterate.
\end{enumerate}
If we define
\[
\Delta_k:=\|V^{(k+1)}-V^{(k)}\|,
\]
then the contraction property gives
\[
\Delta_k\le \beta\Delta_{k-1}\le \beta^k\Delta_0.
\]
So successive differences shrink geometrically, the iteration sequence is Cauchy, and the limit is the unique value-function triple $V^*$. This is the bridge from the fixed-point theorem to an implementable numerical procedure.

\paragraph{Step 6: recover policy objects only after the value functions have been solved.}
Only after value iteration has converged, or after the theoretical fixed point $V^*$ has been established, should we compute
\[
\Gamma_s(z;M)=\arg\max_{a\in\mathcal A}\mathcal W_{sa}(z;M;V^*)
\]
and
\[
\mathcal Z_{sj}(M)=\{z\in Z:j\in\Gamma_s(z;M)\}.
\]
The order matters because the policy objects are derived from the solved value functions; they are not primitive objects that can be specified first. More importantly for the next phase, these measurable choice regions are exactly the objects from which Phase 4 constructs the induced transition matrix and the stationary-distribution block. The endpoint of Phase 3 is therefore not just the statement that $V^*$ exists, but the production of a solved value-function triple together with the measurable policy objects implied by that solution.

\paragraph{Formal summary.}
In this benchmark, the Bellman/value-function solution sequence is:
\begin{align*}
\text{first write }W_{sj}
&\Longrightarrow
\text{then write }V_s=\max_j W_{sj}\\
&\Longrightarrow
\text{rewrite the problem as }V=\mathcal T_MV\\
&\Longrightarrow
\text{use Banach to obtain the unique fixed point}\\
&\Longrightarrow
\text{use value iteration to compute }V^*\\
&\Longrightarrow
\text{recover the policy correspondence and choice regions at the end.}
\end{align*}
Those policy correspondences and choice regions then serve as the direct inputs into Phase 4's transition-matrix and stationary-distribution analysis.

\subsection{Numerical solution of the Bellman block (full procedure)}

The fixed-point theorem tells us that the value functions exist and are unique. A first-year PhD student still needs to see how the fixed point is actually computed. Because $\mathcal T_M$ is a contraction, the constructive solution method is value iteration.

\paragraph{Step N1: initialize.}
Choose grid points or quadrature nodes
\[
z_1,\dots,z_N\in Z
\]
and construct numerical approximations of the continuation kernels. For each active type $j\in\{A,B,C\}$, let $\mathbf Q_j=(q_{j,nm})_{n,m=1}^N$ be a row-stochastic matrix:
\[
q_{j,nm}\ge 0,
\qquad
\sum_{m=1}^N q_{j,nm}=1
\qquad\text{for each }n.
\]
The row-stochastic property is the discrete analog of the statement that $Q_j(\cdot\mid z,M)$ is a probability measure.

Then choose any bounded initial guess $V^{(0)}=(V_A^{(0)},V_B^{(0)},V_C^{(0)})$ on the state grid. A convenient benchmark initialization is
\[
V_s^{(0)}(z_n)=0\quad\forall s\in\{A,B,C\},\ \forall n.
\]
Any other bounded continuous initial guess is also valid, because contraction guarantees convergence from any starting point.

\paragraph{Step N2: continuation integration.}
For each $j\in\{A,B,C\}$ and each grid node $z_n$, compute
\[
\mathcal C_j^{(k)}(z_n):=\sum_{m=1}^N q_{j,nm}V_j^{(k)}(z_m).
\]
This is the discrete approximation to
\[
\int_Z V_j^{(k)}(z')Q_j(dz'\mid z_n,M).
\]
At this step, nothing has been maximized yet. We are only evaluating the expected continuation value implied by the current iterate.

\paragraph{Step N3: Bellman update.}
For each inherited type $s$ and each grid node $z_n$, compute the four action values:
\begin{align*}
\mathcal W_{sA}^{(k)}(z_n)&=\Pi_A(z_n;w,M)-\kappa_{sA}(M)+\beta \mathcal C_A^{(k)}(z_n),\\
\mathcal W_{sB}^{(k)}(z_n)&=\Pi_B(z_n;p_m,M)-\kappa_{sB}(M)+\beta \mathcal C_B^{(k)}(z_n),\\
\mathcal W_{sC}^{(k)}(z_n)&=\Pi_C(z_n;p_m,w,M)-\kappa_{sC}(M)+\beta \mathcal C_C^{(k)}(z_n),\\
\mathcal W_{s,\mathrm{exit}}^{(k)}(z_n)&=0.
\end{align*}
Then update the value function by
\[
V_s^{(k+1)}(z_n)=\max\bigl\{\mathcal W_{sA}^{(k)}(z_n),\mathcal W_{sB}^{(k)}(z_n),\mathcal W_{sC}^{(k)}(z_n),0\bigr\}.
\]
This is the discrete analog of applying the Bellman operator once.

\paragraph{Step N4: successive-difference contraction.}
Define the observable residual
\[
\Delta_k:=\|V^{(k+1)}-V^{(k)}\|.
\]
Because $V^{(k+1)}=\mathcal T_MV^{(k)}$ and $V^{(k)}=\mathcal T_MV^{(k-1)}$, the contraction theorem implies
\[
\Delta_k
=
\|\mathcal T_MV^{(k)}-\mathcal T_MV^{(k-1)}\|
\le \beta\|V^{(k)}-V^{(k-1)}\|
=\beta\Delta_{k-1}.
\]
Applying this recursively yields
\[
\Delta_k\le \beta^k\Delta_0.
\]
So the gap between successive iterates shrinks geometrically.

\paragraph{Step N5: Cauchy bound and convergence.}
For integers $m>n$, the triangle inequality gives
\begin{align*}
\|V^{(m)}-V^{(n)}\|
&\le \sum_{r=n}^{m-1}\|V^{(r+1)}-V^{(r)}\|\\
&= \sum_{r=n}^{m-1}\Delta_r\\
&\le \sum_{r=n}^{m-1}\beta^r\Delta_0\\
&\le \frac{\beta^n}{1-\beta}\Delta_0.
\end{align*}
Because $\beta\in(0,1)$, the right-hand side converges to zero as $n\to\infty$. Therefore the iteration sequence is Cauchy and converges to the unique fixed point $V^*$.

\paragraph{Step N6: a-posteriori error bound and stopping rule.}
The practically useful error bound is
\[
\|V^{(k)}-V^*\|
\le
\sum_{r=k}^{\infty}\|V^{(r+1)}-V^{(r)}\|
\le
\sum_{m=0}^{\infty}\beta^m\Delta_k
=
\frac{1}{1-\beta}\Delta_k.
\]
So a valid stopping rule is: stop iterating once
\[
\frac{\Delta_k}{1-\beta}<\varepsilon_{\mathrm{fp}}.
\]
This rule is operational because the right-hand side is observable during the iteration.

\paragraph{Step N7: recover benchmark objects after convergence.}
Once $V^{(k)}$ has converged to $V^*$, recover the benchmark choice-specific values by substituting $V^*$ into the Bellman indices. Then compute the benchmark argmax correspondence and the benchmark choice regions. In symbols,
\[
\Gamma_s(z;M)=\arg\max_{a\in\mathcal A}\mathcal W_{sa}(z;M;V^*),
\qquad
\mathcal Z_{sj}(M)=\{z\in Z:j\in\Gamma_s(z;M)\}.
\]
Only after the value-function iteration has converged should these final policy objects be recorded for the next phase.

\paragraph{Step N8: transition interface.}
Given stationary measures $\mu_s$, Phase 4 uses
\[
P_{sj}(M)=\frac{\mu_s(\mathcal Z_{sj}(M))}{\mu_s(Z)}.
\]
Because $\mathcal Z_{sj}(M)$ are measurable, $P_{sj}(M)$ is well defined whenever $\mu_s(Z)>0$.

\subsection{What did we just do, and why does it matter?}

We started from the benchmark's economic statement of the incumbent problem and converted it into a mathematical object that can actually be analyzed. The raw Bellman equations are recursive equations in unknown functions. Writing those equations down is only the first step. To know that the model is mathematically well posed, we must know that the equations have a solution, that the solution is unique, and that the solution can be approximated by a convergent algorithm. That is why we introduced the function space $\mathcal V$, the norm $\|\cdot\|$, and the Bellman operator $\mathcal T_M$.

Why did we need a Bellman operator rather than just the Bellman equations themselves? Because Banach's Fixed Point Theorem applies to maps from a metric space into itself. The Bellman equations are a recursive statement equating today's value with today's flow payoff plus tomorrow's continuation value. The Bellman operator is the device that turns that verbal recursion into a self-map on a space of candidate value functions. Once we have a self-map, we can ask mathematically precise questions: does it stay inside the chosen function space, and does it shrink distances between candidate value functions?

Why do we care about contraction? Contraction is the property that makes the recursion stable. Suppose two researchers start from two different guesses for the value functions. If one Bellman update shrinks the distance between those guesses by at least the factor $\beta$, and if $\beta<1$, then repeated Bellman updates force the two sequences of guesses to move toward one another. This means the recursive problem does not wander among multiple incompatible candidate values. Instead it pulls every initial guess toward a single fixed point. That is the mathematical reason value iteration works in this block.

Why does uniqueness of $V^*$ matter for the economics that follow? Every later dynamic or equilibrium object is built on top of the solved value functions. Entry values depend on the value functions. Organizational choice regions depend on comparing choice-specific values, and those values contain continuation values. Stationary distributions depend on the induced policy behavior and the transition structure implied by those policy objects. Market-clearing prices in the stationary equilibrium block must eventually be consistent with those induced objects. If the Bellman system had multiple fixed points, then the downstream objects could inherit that multiplicity. Uniqueness eliminates that ambiguity conditional on fixed prices.

The policy discussion matters for a separate reason. The fixed-point theorem gives us a unique value-function triple, but it does not automatically give us a unique action at every state. If ties occur, then the benchmark policy object is set-valued even though the value functions are unique. That distinction is important. The value functions can be perfectly well defined even when the induced policy correspondence is multi-valued. If a later phase needs a single-valued policy or a row-stochastic transition matrix, tie handling must be added openly rather than hidden inside the notation.

This is exactly why Phase 3 ends where it does. The output of the present phase is not merely the sentence ``the Bellman system has a solution.'' The output is the unique fixed-price value-function triple together with benchmark-faithful measurable policy objects and a transparent statement of what additional convention is needed if later numerical implementation requires a selector. Those are the objects Phase 4 will use when constructing transition matrices and stationary distributions.

\subsection{End-of-Section Recap}

\paragraph{What has been established mathematically.}
For fixed $(w,p_m,M)$, the benchmark dynamic organizational-choice problem defines a Bellman operator on $\mathcal V=\mathcal C_b(Z)^3$. Under the stated regularity assumptions, that operator is a contraction, so the Bellman system has a unique value-function triple. The benchmark argmax correspondence and benchmark choice regions are well defined, and value iteration converges to the unique fixed point with an explicit computable error bound.

\paragraph{What assumptions were used.}
We used the exact benchmark payoff objects $\Pi_A(z;w,M)$, $\Pi_B(z;p_m,M)$, $\Pi_C(z;p_m,w,M)$, the switching-cost matrix $\kappa_{sj}(M)$, the continuation kernels $Q_j(\cdot\mid z,M)$, and discounting with $\beta\in(0,1)$. For the proof, we additionally required bounded continuity of the payoff objects, Feller continuity of the kernels, finiteness of the switching costs, and the compact metric-space structure of $Z$ inherited from Phase 1.

\paragraph{What remains to be derived next.}
The next phase must use the unique fixed-price value functions and the benchmark choice regions to construct the endogenous transition matrix, the stationary distributions across organizational forms, the entry block, and the stationary equilibrium conditions that jointly determine $(w,p_m)$ and the long-run allocation.
